\cotefile{8}
\scanpage{79}
Si p.ex. \deleted{\unsure{$n=p$ est un nb premier}} \inserted{si $\Lambda$ est un corps, (p.ex. $\Lambda=\deleted{\unsure{\mathbb Z/p\mathbb Z}}\mathbb F_p$)}, le $\mathbb L$ sur le dernier $\otimes$ est inutile, on peut écrire
\[
\ell_{\cdot}^{\unsure{G|C}}\simeq\bigl(M\overset{\mathbb L}{\otimes}_{\mathbb Z}\deleted{\unsure{\Lambda}}\mathbb F_p\bigr)\otimes_{\mathbb F_p}\underline{O}_S
\quad .\tag*{$(1.4')$}
\]

\deleted{\unsure{2.|On}}\inserted{1.5.} Supposons qu'on ait un \unsure{revètement|revêtement} \deleted{\illegible{1}}principal $S'$ de $S$ de groupe \deleted{\unsure{ordinaire}} $\pi$ discret ou profini, et que $M$ provienne par descente d'une \deleted{\unsure{revêtement}} représentation de $\pi$ dans un groupe \unsure{abélien|abèlien} $M_0$ (qu'on suppose continue pour $\pi$ profini)\,; p.ex. si $M$ est isotrivial (p.ex. à fibres de t. fini et $S$ \deleted{connexe} normal) et si $S$ connexe, alors on sera dans la situation dite, avec $S'$ le rev. universel de $S$.... On a alors un \unsure{hom|hom.} de topos
\diagram{\begin{tikzcd}[ampersand replacement=\&,column sep=large,row sep=0pt]
(1.6) \& \bigl(S\deleted{\unsure{,\underline{O}_S}}\bigr) \arrow[r,"f"] \& \bigl(B_{\pi}\deleted{\unsure{,\mathbb Z|,\Lambda}}\bigr) \\
\& \inserted{\substack{\text{(top. étale,}\\\text{ou fppf}\\\text{ou \unsure{fpqc|fpqe}}}} \&
\end{tikzcd}}
avec
\[
M=f^{\deleted{*}-1}(M_0)\tag*{(1.7)}
\]

On aura alors, si on \unsure{annèle|annéle} $S$ et $B_{\pi}$ par l'anneau constant $\Lambda$, et envisage $f$ en conséquence
\[
M\overset{\mathbb L}{\otimes}_{\mathbb Z}\deleted{\unsure{\mathbb F_p}}\Lambda
=f^*\bigl(M_0\overset{\mathbb L}{\otimes}_{\mathbb Z}\Lambda\bigr)\tag*{\deleted{\unsure{(1.7)}}}
\]
et par suite on peut écrire
\[
\ell_{\cdot}^{G}\simeq f^*\bigl(M_0\overset{\mathbb L}{\otimes}_{\mathbb Z}\Lambda\bigr)\overset{\mathbb L}{\otimes}_{\Lambda}\underline{O}_S
\quad ,\tag*{(1.8)}
\]
où le $\mathbb L$ sur $\otimes_{\Lambda}$ est inutile si $\Lambda$ est un corps....
% archive sheet number: 78
\scanpage{80}
2. Notons que \deleted{si $\underline{O}_S$ est de car. $\unsure{p>0}$,} l'on tire de (1.3)
\[
\left\{\begin{aligned}
\omega_{G}&=\underline{M}\otimes_{\mathbb Z}\underline{O}_S\\
\underline{n}_G&=\underline{\operatorname{Tor}}^{\mathbb Z}_1(M,\underline{O}_S)
\end{aligned}\right.\tag*{(2.1)}
\]
ce qui s'écrit, si $\deleted{\underline{O}_S}S$ est de car. $p>0$
\[
\left\{\begin{aligned}
\omega_G&=(M\otimes_{\mathbb Z}\unsure{F|\mathbb F}_p)\otimes_{\mathbb F_p}\underline{O}_S=M_p\otimes_{\mathbb F_p}\underline{O}_S\\
\underline{n}_G&\simeq {}_pM\otimes_{\mathbb F_p}\underline{O}_S
\end{aligned}\right.\tag*{(2.\unsure{1|2})}
\]
et si $M$ est annulé par $p$, i.e. $G$ annulé par $p$, cela devient
\[
\omega_G\simeq M\otimes_{\mathbb F_p}\underline{O}_S
\quad ,\qquad \underline{n}_G\simeq M\otimes_{\mathbb F_p}\underline{O}_S
\tag*{$(2.2')$}
\]

On a d'autre part une classe canonique
\[
\xi\in\operatorname{Ext}_{\underline{O}_S}^{2}\bigl(M\otimes_{\mathbb Z}\underline{O}_S,\underline{\operatorname{Tor}}^{\mathbb Z}_1(M,\underline{O}_S)\bigr)
\]
qu'on peut \unsure{interpreter|interpréter}, en présence d'un $\Lambda$, comme provenant d'un
\[
\xi_{\Lambda}\in\operatorname{Ext}_{\Lambda}^{2}\bigl(M\otimes_{\mathbb Z}\Lambda,\operatorname{Tor}^{\mathbb Z}_1(M,\Lambda)\bigr)\quad ;
\]

Supposons p.ex. que $\Lambda=\mathbb Z/n\mathbb Z$, et que $\deleted{M\otimes_{\mathbb Z}}\inserted{M}$ soit un $\Lambda$-module loc. libre de t.f. Alors il en est de $\unsure{\overline{m}|\eqsim}$ de $M\otimes_{\mathbb Z}\Lambda\simeq M$ et de $\operatorname{Tor}^{\mathbb Z}_1(M,\Lambda)\simeq {}_nM\simeq M$, et
\[
\xi_{\Lambda}\in H^2(S,\operatorname{End}(M))\quad ,
\]
et \deleted{si $\deleted{\xi\in}\inserted{S}\deleted{\underline{O}_S}$ est plat sur $\Lambda$ (p.ex. \unsure{$n=p$|en car. $p$})}
\[
\xi\in H^2(S,\operatorname{End}(\deleted{M}E))\qquad (E=M\otimes_{\Lambda}\underline{O}_S)
\]
et $\xi$ est déduit de $\xi_{\Lambda}$ par l'hom can.
% archive sheet number: 79
\scanpage{81}
% Faint hooked trace at the upper left edge of the sheet.
\[
H^*(S,\operatorname{End}(M))\mathrel{\unsure{\relbar|\longrightarrow}}
H^*(S,\operatorname{End}(M)\otimes_{\Lambda}\underline{O}_S)\simeq H^{\unsure{2|*}}(S,\operatorname{End}(E))
\]

Enfin, sous \unsure{les conditions|la condition} de \unsure{1.5|(1.5}, $M_0$ \unsure{etant|étant} aussi plat sur $\Lambda$, on aura que $\unsure{\xi|\zeta}_{\Lambda}$ est l'image d'un élément
\[
\xi_0\in H^2(B_{\pi},\operatorname{End}(M_0))\quad .
\]

(NB Pour calculer $\xi_0$, par fonctorialité en $\pi$, on est ramené aussitôt au cas ``universel'' $\pi=\operatorname{Aut}(M_0)$, et on peut \unsure{évidemment|évidement} supposer $M_0=\Lambda^d=(\mathbb Z/n\mathbb Z)^d$, donc $\pi=\operatorname{Gl}(d,\mathbb Z/n\mathbb Z)$....) D'après \uline{Illusie}, c'est l'obstruction à remonter la représentation de $\pi$ dans $M_0$ en une représentation dans un module libre sur $\mathbb Z/\unsure{n|p}^2\mathbb Z$. \inserted{Lorsque $n=p$ (premier)} Elle est \unsure{évidemment|évidement} nulle si $M_0$ est de rang $\leq 1$, mais sans doute pas nulle si $M_0$ de rang $d\geq 2$, dans le cas universel $\pi=\deleted{\unsure{A}}\operatorname{Gl}(d,\mathbb F_p)$.

Notons que si le rev. $S'$ de $S$ est le rev. universel, alors
\[
H^2(B_{\pi},-)\longrightarrow H^2(S,-)
\]
est \uline{injectif}, donc on \unsure{trouvera|trouve} un élément non nul \deleted{dans} $\xi_{\Lambda}\in H^2(S,\operatorname{End}(M))$.
% archive sheet number: 80
